\chapter{Na\"ive set theory}

This chapter serves as a foundation, although not a very strong one, for mathematical analysis. It
is sufficient to persue the journey of analysis, but it is in no means rigorous or complete. We
would not cover the ordinals, cardinals, and the various, almost paradoxical consequences of some of
its axioms. We would also not go over proving that some objects are sets. We would not abstract so
much that everything in the universe are sets either. That's why this chapter title is called
Na\"ive set theory. To the readers that want a rigorous development of set theory, refer to any the
vast repertoire of textbooks in the field. One of my favorite being
\citeauthor{enderton-set-theory}'s \citetitle{enderton-set-theory} \cite{enderton-set-theory}, and
\citeauthor{halmos-set-theory}'s \citetitle{halmos-set-theory} \cite{halmos-set-theory} which is
albeit, not as Na\"ive as the title would've suggested. \citeauthor{goldrei-1996}'s
\citetitle{goldrei-1996} is also a good book, although it leans more towards analysis side of
things.

\section{Sets and collections}

The vague notion of set is captured in the following definition:

\begin{df}{A na\"ive definition of sets}{}
  A \textbf{set} is a \emph{collection} of object that contain things. If an object \(a\) is in a
  set \(A\), i.e., it is a member of \(A\), then we write \(a \in A\). If it isn't a member of
  \(A\), then \(a \notin A\).
\end{df}

\paragraph{Collection}

is a vague, non-rigorous term that describes a group of things. It's actually more a philosophical
matter than a mathematical matter. At the tip of the iceberg, a collection is literally just a group
of things. They have no structure, and no relations to each other. In a sense, it's just a choice of
symbols and has no other functions. Set theory is then overlaid on top of the notion of collections.

A set is just a collection that follows a list of rules, specifically the Zermelo-Frankel set
theory. But I don't want to get in too deep. So we shall roll with it and pretend that a collection
is just a group of things like two apples in a basket, people on Earth, chairs that are on a table,
table that are on a chair, etc.

However, a collection is different from sets in its structures. As you'll see later, the structure
of sets builds on top of collections so that two sets can interact. We can rigorously compare,
combine, destroy, or even nest their elements. But, a collection can't as it is just a vague
description of groups.
